% figures/foundations/state_hierarchy.tex

\begin{figure}[ht]
\centering
\begin{tikzpicture}[
  node distance=0.85cm,
  box/.style={
    draw,
    rounded corners,
    align=center,
    minimum width=4.8cm,
    minimum height=0.7cm
  },
  arrow/.style={->, thick}
]
\node[box] (kappa) {Curvature Law\\$\kappa(s)$};
\node[box, below=of kappa] (pose2) {pose2\\intrinsic planar state};
\node[box, below=of pose2] (pose3) {pose3\\railway runtime state};
\node[box, below=of pose3] (runtime) {Operational Runtime State\\$x(s)$};
\node[box, below=of runtime] (vehicle) {Vehicle State\\$x_{\mathrm{vehicle}}$};
\node[box, below=of vehicle] (cg) {CG State\\$x_{\mathrm{CG}}$};
\node[box, below=of cg] (opt) {Optimization State\\$x_{\mathrm{opt}}$};

\draw[arrow] (kappa) -- (pose2);
\draw[arrow] (pose2) -- (pose3);
\draw[arrow] (pose3) -- (runtime);
\draw[arrow] (runtime) -- (vehicle);
\draw[arrow] (vehicle) -- (cg);
\draw[arrow] (cg) -- (opt);
\end{tikzpicture}

\caption{AIM state hierarchy from intrinsic curvature to optimization state.}
\label{fig:state-hierarchy}
\end{figure}
